
This paper reports three empirical measurements relevant to orbit-based positional encodings for weight space learning.

First, the input/output channel permutation group is a functionally exact symmetry for Conv2d, Linear, and MultiheadAttention layers in both CNN Zoo and Transformer Zoo checkpoints (|Delta acc| = 0.000000 across 4,500 runs). This result is stronger than required by the pre-specified threshold.

Second, orbit-PE is computable for all layer types with 1.167x mean overhead relative to sequential PE, using a unified codebase with no architecture-conditional branches. Overhead is consistent across layer types (Conv2d: 1.168x, Linear: 1.168x, MultiheadAttention: 1.126x).

Third, variance decomposition across 1,000 CNN Zoo models and 50 training epochs reveals a layer-type stratification: the permutation-orbit variance fraction is 0.637 for Conv2d layers and 0.133 for Linear (FC) layers, with an overall ratio of 0.3479 — below the 0.60 gate threshold. This gate failure prevented execution of the planned cross-architecture training experiments (H-M3, H-C1). An ancillary finding is that the permutation variance ratio decreases during training (approximately 0.49 to 0.28 over 50 epochs).

These results collectively indicate that permutation orbit encoding alone is insufficient as the basis for cross-architecture WSL due to GL-orbit dominance in fully-connected and attention layers. The findings motivate a hybrid positional encoding (permutation orbit-PE for Conv2d + GL-invariant trace features for Linear/attention) as the next iteration. Whether such a hybrid achieves the originally targeted tau_retention >= 0.70 (relaxed to >= 0.65 in light of the H-M2 finding) remains untested. The variance decomposition infrastructure developed here is directly applicable to that next step.

